\subsection{Agentic Hybrid System：组件、流程、资产与新增攻击面}

slides 10--18 从 stand-alone LLM 转到 compound system。Agentic system 通常由用户、host、一个或多个模型、数据源、memory、tools 与外部环境组成。一次请求经过 prompt 构造、模型决策、工具执行、结果回填和最终响应。安全目标仍包括 confidentiality、integrity、availability，但新增了模型、系统提示、memory、tool credentials、action history 与外部资源等保护对象；LLM 的不确定性又扩大了传统系统的输入边界。

\lecturefigure{slides-images/slide-010.png}{官方 slide 10：LLM Safety 与 LLM Agent Safety 的范围差异}
\lecturefigure{slides-images/slide-011.png}{官方 slide 11：LLM Agent 与 Agentic System 的定义}
\lecturefigure{slides-images/slide-012.png}{官方 slide 12：Agentic System 作为 hybrid/compound system}
\lecturefigure{slides-images/slide-013.png}{官方 slide 13：Agentic Hybrid System 的端到端使用流程}
\lecturefigure{slides-images/slide-014.png}{官方 slide 14：Hybrid System 的 security 与 safety goals}
\lecturefigure{slides-images/slide-015.png}{官方 slide 15：相对传统系统新增的保护目标}
\lecturefigure{slides-images/slide-016.png}{官方 slide 16：使用 LLM 带来的 attack surface 扩张之一}
\lecturefigure{slides-images/slide-017.png}{官方 slide 17：Attack surface 扩张之二}
\lecturefigure{slides-images/slide-018.png}{官方 slide 18：Attack surface 扩张的完整状态}

\begin{knowledgebox}{术语消化：Agent 安全的四类资产}
\begin{tabular}{p{2.7cm}p{4.2cm}p{5.3cm}}
\toprule
\textbf{资产} & \textbf{示例} & \textbf{失守后果} \\
\midrule
Secrets & system prompt、token、API key & 泄露权限与内部规则 \\
State & memory、计划、会话与业务状态 & 持久污染、错误决策、跨会话传播 \\
Actions & tool call、数据库写入、外部消息 & 越权、破坏、欺诈和不可逆副作用 \\
Trust evidence & logs、provenance、policy decision & 无法审计、归因或证明修复有效 \\
\bottomrule
\end{tabular}
\end{knowledgebox}

\begin{importantbox}{系统边界决定安全边界}
模型本身“没有泄露 secret”并不等于系统安全：模型可能生成查询语句、shell command 或 URL，随后由高权限组件执行。威胁建模必须沿数据流和控制流追踪每个 trust boundary，明确谁能读、谁能写、谁能批准、谁能回滚。
\end{importantbox}

